\section*{Chapitre \ref{ch:ll:protocols-ii-binary-exchange-protocols} --- Protocoles II~: protocoles binaires d'échange}

\begin{solution}{exo:ll:protocols-ii-binary-exchange-protocols:1}
Le premier paquet porte 1\,000--1\,002, donc le numéro suivant attendu est 1\,003~: les messages 1\,003 à 1\,006 manquent (un
\omterm{def:ll:protocols-ii-binary-exchange-protocols:multicast}{trou de séquence} de quatre). Le numéro du \omterm{def:ll:protocols-i-fix:session}{heartbeat} est le prochain attendu, donc le paquet à 1\,007 portait trois messages (1\,007--1\,009) et
rien d'autre n'a été perdu après lui.
\end{solution}

\begin{solution}{exo:ll:protocols-ii-binary-exchange-protocols:2}
$\mathtt{0x000F42A4} = 1\,000\,100$ dix-millièmes~: 100.01.
\end{solution}

\begin{solution}{exo:ll:protocols-ii-binary-exchange-protocols:3}
$0.002^2 = 4 \times 10^{-6}$ avec des lignes indépendantes. Avec le commutateur partagé, $10^{-4} + (1 - 10^{-4}) \times 4 \times 10^{-6}
\approx 1.04 \times 10^{-4}$~: le commutateur seul explique 96~\% des pertes.
\end{solution}

\begin{solution}{exo:ll:protocols-ii-binary-exchange-protocols:4}
La probabilité stationnaire de l'état dégradé est $0.0005/0.1005 \approx 0.50~\%$, donc la ligne perd $0.8 \times 0.50~\% \approx 0.40~\%$
des paquets~; une période dégradée dure $1/0.1 = 10$ paquets en moyenne. Les pertes arrivent donc par salves de plusieurs paquets à la fois~: une
demande de retransmission les couvre, mais au moment où elles se produisent, souvent en pleine rafale de trafic, et une seule demande peut
ne pas suffire.
\end{solution}

\begin{solution}{exo:ll:protocols-ii-binary-exchange-protocols:5}
Un paquet \omterm{def:ll:protocols-ii-binary-exchange-protocols:multicast}{multicast} atteint tous les abonnés à la fois, sans connexion, accusé de réception ni retransmission propre à chaque abonné pour
ralentir l'éditeur, et aucun abonné ne peut retenir les autres. Un abonné renonce à une livraison fiable et ordonnée~: il doit
détecter les trous à partir des numéros de séquence, arbitrer les lignes redondantes et se rétablir seul par retransmission ou instantanés.
\end{solution}

\begin{solution}{exo:ll:protocols-ii-binary-exchange-protocols:6}
Le message devient 37 octets de champs plus un en-tête de 8 octets, soit 45 contre 36, soit 25~\% de bande passante et de cache en plus pour les Add Orders.
Ce qu'il économise, c'est l'assemblage d'un entier de six octets~; la mesure montre que ce gain n'est pas mesurable, tandis que
l'en-tête apporte un numéro de version et une longueur qui permettent aux anciens lecteurs de sauter les champs qu'ils ne connaissent pas.
\end{solution}

\begin{solution}{exo:ll:protocols-ii-binary-exchange-protocols:7}
Ajouter \texttt{soup\_client} et \texttt{soup\_server} aux protocoles que parcourt le générateur, et faire générer pour les messages
de longueur variable (\texttt{U}, \texttt{S}, \texttt{+}) un lecteur qui renvoie la charge utile sous forme de vue~; la trame porte une longueur u16 big-endian
qui compte l'octet de type. Le test lit \texttt{fixture\_soup.bin} trame par trame et compare chacune avec
\texttt{fixture\_soup.csv}.
\end{solution}

\begin{solution}{exo:ll:protocols-ii-binary-exchange-protocols:8}
Il n'arbitre pas~: pendant qu'il lit une ligne, les paquets de l'autre sont ignorés, si bien que chaque commutateur perd ce qui
est arrivé en premier sur l'autre ligne, et un trou sur la ligne A n'est comblé que si les copies de la ligne B sont encore à venir. Le gestionnaire doit lire les deux lignes
en permanence et garder, pour chaque numéro de séquence, la première copie reçue de l'une ou de l'autre.
\end{solution}

\begin{solution}{pb:ll:protocols-ii-binary-exchange-protocols:1}
\begin{enumerate}
\item $50\,000 \times 23\,400 = 1.17 \times 10^9$ paquets~; une ligne en perd environ 117\,000.
\item $10^{-8}$, soit environ 11.7 pertes par session.
\item Parce que les pertes n'étaient pas indépendantes~: les deux lignes ont été coupées ensemble pendant 25 millisecondes (17 messages), et la ligne B a perdu un
paquet pendant que la ligne A était hors service.
\item Un aller-retour de retransmission, \qty{0.5}{\milli\second}, pendant lequel le carnet ne peut pas être tenu pour fiable.
\item $10^{-6} + (1 - 10^{-6}) \times 10^{-8} \approx 1.01 \times 10^{-6}$~: environ 1\,182 pertes par session, cent fois plus.
\item Environ 0.40~\% par ligne, et $1.6 \times 10^{-5}$ après arbitrage si les lignes sont indépendantes.
\item Les pertes arrivent par séries aux moments les plus chargés~; la demande, la réponse et le rejeu ont lieu pendant que le marché bouge, et une
série plus longue que la limite de retransmission exige un instantané.
\item Un commutateur partagé, une carte réseau ou un serveur partagé, un cross-connect ou un opérateur partagé, une alimentation commune~; il ne peut pas supprimer les
pannes propres à l'éditeur, qui sont communes aux deux lignes.
\item Quand l'écart dépasse ce que le serveur conserve ou sert (1\,000 messages par demande, une fenêtre de 100\,000), quand le
serveur est injoignable, ou quand le gestionnaire démarre en retard.
\item En moyenne, une demi-période plus l'instantané, soit $0.5 + 0.02 = \qty{0.52}{\second}$~; au pire, une période complète,
\qty{1.02}{\second}.
\item Environ 26\,000 messages en moyenne (51\,000 au pire), appliqués en environ \qty{0.21}{\milli\second} (\qty{0.41}{\milli\second}).
\item Ceux que l'instantané reflète déjà~: les numéros de séquence jusqu'à celui indiqué par le marqueur de début de l'instantané
(LastMsgSeqNumProcessed du CME).
\item \textbf{Résultat nommé.} Avec des pertes indépendantes de $10^{-4}$ par ligne, le flux arbitré perd $10^{-8}$ par paquet,
soit environ 12 paquets par session à 50\,000 par seconde~; une cause commune de $10^{-6}$ le porte à $1.01 \times 10^{-6}$, soit environ 1\,180 par
session. La reprise sur un cycle d'instantanés d'une seconde prend \qty{0.52}{\second} en moyenne et \qty{1.02}{\second} au pire,
contre \qty{0.5}{\milli\second} pour une retransmission.
\item Marquer le carnet comme périmé, cesser de coter dessus (ou ne coter que de façon défensive) et annuler ce qui en dépend jusqu'à sa reconstruction.
\item Cela ramène l'attente moyenne à \qty{50}{\milli\second}, au prix de dix fois la bande passante des instantanés pour
chaque abonné, dont la plupart n'en ont jamais besoin.
\item Attendre la ligne la plus lente ajoute son délai à chaque message~; l'intérêt de deux lignes est de prendre la copie la plus rapide.
\item Rejouer les lignes enregistrées avec des pertes injectées, une ligne à la fois puis les deux ensemble (comme le fait le simulateur), et vérifier que
le carnet reconstruit égale celui d'une exécution sans perte.
\item La ligne, la plage de séquence, l'heure, si le paquet a été récupéré et comment, et combien de temps le carnet est resté périmé.
\item Presque nulle part dans le décodage~: quelques nanosecondes par message~; la latence se trouve dans le réseau, le noyau et le carnet.
\item Une protection contre les pannes indépendantes d'un chemin, et non contre les pannes communes aux deux.
\end{enumerate}
\end{solution}

\begin{solution}{iq:ll:protocols-ii-binary-exchange-protocols:1}
Un paquet atteint tous les abonnés à la fois, et l'éditeur n'attend personne. Chaque paquet porte le numéro de séquence de son premier
message et un compte~; un paquet qui commence au-dessus du numéro suivant attendu révèle un trou, et les \omterm{def:ll:protocols-i-fix:session}{heartbeats} portent le numéro
suivant, de sorte que les pertes soient détectées même quand rien ne se négocie.
\iqlookfor{équité et passage à l'échelle~; numéros de séquence et \omterm{def:ll:protocols-i-fix:session}{heartbeats}.}
\end{solution}

\begin{solution}{iq:ll:protocols-ii-binary-exchange-protocols:2}
Conserver le prochain numéro de séquence attendu. Pour chaque paquet venant de l'une ou l'autre ligne~: s'il se termine avant ce numéro, le rejeter~; s'il commence
après lui, consigner un trou (perdu sur les deux lignes) et lancer la reprise~; sinon livrer les nouveaux messages et avancer. Lire les deux lignes
sans attendre l'une ou l'autre~; mettre en tampon ce qui arrive pendant la reprise.
\iqlookfor{un seul compteur, doublons écartés, trous comblés.}
\end{solution}

\begin{solution}{iq:ll:protocols-ii-binary-exchange-protocols:3}
Un objet ne contenant qu'un pointeur vers le message, dont les accesseurs lisent chaque champ à son décalage fixe lorsqu'on le demande. Rien n'est
copié, rien n'est alloué, les champs non lus ne coûtent rien, et il n'y a aucun appel indirect~: le coût mesuré est la moitié de celui
d'un décodeur qui copie dans une structure et appelle un rappel.
\iqlookfor{vues sur des octets et accès paresseux aux champs.}
\end{solution}

\begin{solution}{iq:ll:protocols-ii-binary-exchange-protocols:4}
S'abonner aux canaux incrémentaux et mettre en tampon~; attendre un instantané de chaque instrument~; le charger~; supprimer les messages tamponnés jusqu'au
numéro de séquence que reflète l'instantané~; appliquer le reste dans l'ordre~; puis traiter en temps réel. Vérifier avec la somme de contrôle de l'instantané
lorsqu'il y en a une.
\iqlookfor{tampon, instantané, alignement sur le numéro de séquence.}
\end{solution}

\begin{solution}{iq:ll:protocols-ii-binary-exchange-protocols:5}
Guère~: un échange d'octets ne coûte qu'une instruction, et la mesure du chapitre ne montre aucune différence entre les
flyweights big-endian et little-endian. Ce qui compte, c'est de ne pas copier, de ne pas allouer, de ne pas appeler à travers des pointeurs, et des décalages fixes qui n'ont besoin d'aucune recherche.
\iqlookfor{mesurer~; les coûts sont ailleurs.}
\end{solution}

\begin{solution}{iq:ll:protocols-ii-binary-exchange-protocols:6}
Une cause commune~: des commutateurs, cartes, câbles, opérateurs ou alimentations partagés~; l'éditeur lui-même (les deux lignes perdent la même plage à la
source)~; l'hôte récepteur (les interruptions des deux lignes sur un même cœur surchargé, un tampon de socket plein, un gestionnaire lent). Comparer les
instants et plages des trous entre les lignes et avec les autres abonnés, et vérifier les compteurs de rejets de l'hôte.
\iqlookfor{l'indépendance comme hypothèse à tester, y compris l'hôte récepteur.}
\end{solution}
